\documentclass[11pt]{article}

% ---- packages ----
\usepackage[utf8]{inputenc}
\usepackage[T1]{fontenc}
\usepackage{lmodern}
\usepackage[margin=1in]{geometry}
\usepackage{amsmath,amssymb,amsthm,mathtools}
\usepackage{braket}
\usepackage{enumitem}
\usepackage{xcolor}
\usepackage[colorlinks=true,linkcolor=blue,urlcolor=blue]{hyperref}

% ---- macros (subset of the course preamble) ----
\newcommand{\R}{\mathbb{R}}
\newcommand{\C}{\mathbb{C}}
\newcommand{\E}{\mathbb{E}}
\newcommand{\Prob}{\mathbb{P}}
\newcommand{\ii}{\mathrm{i}}
\newcommand{\ee}{\mathrm{e}}
\DeclareMathOperator{\tr}{tr}
\DeclareMathOperator{\Tr}{Tr}
\DeclareMathOperator{\Var}{Var}
\DeclareMathOperator{\diag}{diag}
\newcommand{\norm}[1]{\left\lVert #1 \right\rVert}
\newcommand{\abs}[1]{\left\lvert #1 \right\rvert}
\newcommand{\bigo}{\mathcal{O}}
\newcommand{\id}{\mathbb{I}}
\newcommand{\dg}{^{\dagger}}
\newcommand{\proj}[1]{\ket{#1}\!\bra{#1}}
\newcommand{\CNOT}{\mathrm{CNOT}}
\newcommand{\QFT}{\mathrm{QFT}}
\newcommand{\defeq}{\vcentcolon=}

\newcommand{\sol}{\smallskip\noindent\textit{Solution.}\ }
\newcommand{\problem}[1]{\bigskip\noindent\textbf{Problem #1.}\ }

\title{\textbf{Quantum Scientific Computing (QSC2026)}\\[0.3em]\large Homework~1 --- Solutions}
\author{}
\date{}

\begin{document}
\maketitle
\vspace{-2.4em}
\noindent\small These solutions show the intended level of detail; many admit
equivalent routes. Numerical answers are exact where a closed form exists and
rounded to four digits otherwise.
\normalsize

% ======================================================================
\problem{1 (Single qubit)}
$\ket{\psi}=a\ket0+b\ket1$ with $a=\cos\tfrac\pi8$, $b=\ee^{\ii\pi/4}\sin\tfrac\pi8$.

\sol
\textbf{(a)} $\abs a^2+\abs b^2=\cos^2\tfrac\pi8+\sin^2\tfrac\pi8=1$. The
computational-basis probabilities are
$\Prob(0)=\cos^2\tfrac\pi8=\tfrac{2+\sqrt2}{4}\approx0.8536$ and
$\Prob(1)=\sin^2\tfrac\pi8=\tfrac{2-\sqrt2}{4}\approx0.1464$.

\textbf{(b)} With $\bar a b=\cos\tfrac\pi8\sin\tfrac\pi8\,\ee^{\ii\pi/4}
=\tfrac12\sin\tfrac\pi4\,\ee^{\ii\pi/4}=\tfrac{\sqrt2}{4}\ee^{\ii\pi/4}$,
\[
 \braket{X}=2\Re(\bar a b)=\tfrac{\sqrt2}{4}\cdot 2\cos\tfrac\pi4=\tfrac12,\quad
 \braket{Y}=2\Im(\bar a b)=\tfrac12,\quad
 \braket{Z}=\abs a^2-\abs b^2=\cos\tfrac\pi4=\tfrac{\sqrt2}{2}.
\]
Bloch vector $\mathbf r=(\tfrac12,\tfrac12,\tfrac{\sqrt2}{2})$, with
$\abs{\mathbf r}^2=\tfrac14+\tfrac14+\tfrac12=1$ (pure state).

\textbf{(c)} $\Prob(+)=\abs{\braket{+|\psi}}^2=\tfrac12\abs{a+b}^2
=\tfrac12\bigl(1+2\Re(\bar a b)\bigr)=\tfrac12(1+\braket{X})=\tfrac12(1+\tfrac12)=\tfrac34$.
(The identity $\Prob(+)=\tfrac12(1+\braket{X})$ is the general rule for an $X$-basis
measurement.)

\textbf{(d)} A $Z$-measurement returns $\pm1$ with variance
$\Var=1-\braket{Z}^2=1-\tfrac12=\tfrac12$; the empirical mean has error
$\approx z\sqrt{\Var/N}$ at confidence level set by $z$. Setting this to
$\epsilon=0.05$ with $z\approx2$ gives $N\approx z^2\Var/\epsilon^2
=4\cdot0.5/0.0025\approx800$. The point is the shot-noise law $N=\Theta(1/\epsilon^2)$
(hundreds to a few thousand shots here).

% ======================================================================
\problem{2 (Separability test)}
A state $a\ket{00}+b\ket{01}+c\ket{10}+d\ket{11}$ factorizes iff $ad=bc$.

\sol
\textbf{(a)} $a=b=c=d=\tfrac12$: $ad=\tfrac14=bc$, so it is a \emph{product} state,
$=\tfrac12(\ket0+\ket1)(\ket0+\ket1)=\ket{+}\otimes\ket{+}$.

\textbf{(b)} $a=\tfrac1{\sqrt2},\,d=-\tfrac1{\sqrt2},\,b=c=0$: $ad=-\tfrac12\neq0=bc$,
so it is \emph{entangled} (the Bell state $\ket{\Phi^-}$).

\textbf{(c)} $a=b=c=\tfrac12,\,d=-\tfrac12$: $ad=-\tfrac14\neq\tfrac14=bc$, so it is
\emph{entangled}.

\textbf{(d)} If $\tfrac1{\sqrt2}(\ket{00}+\ket{11})=(\alpha\ket0+\beta\ket1)
(\gamma\ket0+\delta\ket1)$ then matching coefficients gives
$\alpha\gamma=\beta\delta=\tfrac1{\sqrt2}$ (so $\alpha,\beta,\gamma,\delta\neq0$) but
$\alpha\delta=\beta\gamma=0$ (coefficients of $\ket{01},\ket{10}$), a contradiction.
Hence no factorization exists; equivalently $ad-bc=\tfrac12\neq0$.

% ======================================================================
\problem{3 (Gates, reversibility, no-cloning)}

\sol
\textbf{(a)} $(H\otimes\id)\ket{00}=\ket{+}\ket0=\tfrac1{\sqrt2}(\ket{00}+\ket{10})$;
then $\CNOT$ (control first) sends $\ket{10}\mapsto\ket{11}$, giving
$\tfrac1{\sqrt2}(\ket{00}+\ket{11})=\ket{\Phi^+}$.

\textbf{(b)} Toffoli: $\ket{a,b,c}\mapsto\ket{a,b,\,c\oplus(a\wedge b)}$, so with
$c=0$ it writes $a\wedge b$ onto the third qubit. It is its own inverse, hence
reversible. The bare map $(a,b)\mapsto a\wedge b$ is $2\to1$ and not injective (both
$01$ and $00$ give $0$), so it cannot be a unitary on two qubits; keeping $a,b$ and
writing the result on a fresh qubit restores invertibility.

\textbf{(c)} Take $U=\CNOT$ and $\ket{s}=\ket0$: then $U\ket0\ket0=\ket{00}$ and
$U\ket1\ket0=\ket{11}$ as required. On a general input,
$U(\alpha\ket0+\beta\ket1)\ket0=\alpha\ket{00}+\beta\ket{11}$, whereas a true copy
would be $(\alpha\ket0+\beta\ket1)^{\otimes2}
=\alpha^2\ket{00}+\alpha\beta\ket{01}+\alpha\beta\ket{10}+\beta^2\ket{11}$. These
agree only when $\alpha\beta=0$, i.e.\ for basis inputs; $U$ entangles instead of
copying. Linearity forbids a universal cloner---this is the no-cloning theorem.

% ======================================================================
\problem{4 (Teleportation)}
Message $\ket\psi=\alpha\ket0+\beta\ket1$ on qubit~1; shared
$\tfrac1{\sqrt2}(\ket{00}+\ket{11})$ on qubits~2 (Alice) and~3 (Bob).

\sol
\textbf{(a)} The joint state
$\tfrac1{\sqrt2}(\alpha\ket{000}+\alpha\ket{011}+\beta\ket{100}+\beta\ket{111})$
re-expressed in the Bell basis of qubits $1,2$ is
\[
 \tfrac12\Bigl[\ket{\Phi^+}(\alpha\ket0+\beta\ket1)+\ket{\Phi^-}(\alpha\ket0-\beta\ket1)
 +\ket{\Psi^+}(\alpha\ket1+\beta\ket0)+\ket{\Psi^-}(\alpha\ket1-\beta\ket0)\Bigr]_3 .
\]
With the standard $\CNOT_{1\to2}$ then $H_1$ measurement giving bits $(m_1,m_2)$, the
four outcomes and Bob's pre-correction states are
$(0,0)\!\to\!\ket\psi$, $(0,1)\!\to\! X\ket\psi$, $(1,0)\!\to\! Z\ket\psi$,
$(1,1)\!\to\! XZ\ket\psi$, each with probability $\tfrac14$.

\textbf{(b)} For $(m_1,m_2)=(1,1)$ the pre-correction state is
$XZ\ket\psi=\alpha\ket1-\beta\ket0$. Bob applies $X$ then $Z$ (i.e.\ $Z^{m_1}X^{m_2}$):
$X(\alpha\ket1-\beta\ket0)=\alpha\ket0-\beta\ket1=Z\ket\psi$, then
$Z(\alpha\ket0-\beta\ket1)=\alpha\ket0+\beta\ket1=\ket\psi$. Restored.

\textbf{(c)} Before the bits arrive, all four outcomes are equally likely, so Bob's
density operator is $\tfrac14\sum_{P\in\{\id,X,Z,XZ\}}P\proj{\psi}P\dg=\tfrac12\id$,
the maximally mixed state---independent of $\ket\psi$. Nothing Alice does changes
Bob's local statistics until the two classical bits (sent no faster than light)
arrive; hence no faster-than-light signaling.

% ======================================================================
\problem{5 (Phase kickback; Deutsch; Deutsch--Jozsa)}

\sol
\textbf{(a)} $U_f\ket{x}\tfrac1{\sqrt2}(\ket0-\ket1)
=\tfrac1{\sqrt2}\ket x(\ket{f(x)}-\ket{1\oplus f(x)})
=(-1)^{f(x)}\ket x\tfrac1{\sqrt2}(\ket0-\ket1)=(-1)^{f(x)}\ket x\ket-$, since the pair
$(\ket{f(x)},\ket{1\oplus f(x)})$ equals $(\ket0,\ket1)$ if $f(x)=0$ and
$(\ket1,\ket0)$ if $f(x)=1$.

\textbf{(b)} The four functions: $f\equiv0$ and $f\equiv1$ are \emph{constant};
$f(x)=x$ and $f(x)=1-x$ are \emph{balanced}. Deutsch prepares $\ket+\ket-$; phase
kickback gives the first qubit $\propto(-1)^{f(0)}\ket0+(-1)^{f(1)}\ket1$. Applying
$H$: if $f(0)=f(1)$ the qubit is $\pm\ket+\xrightarrow{H}\pm\ket0$ (measure $0$); if
$f(0)\neq f(1)$ it is $\pm\ket-\xrightarrow{H}\pm\ket1$ (measure $1$). So the single
query returns $0$ iff $f$ is constant.

\textbf{(c)} In Deutsch--Jozsa the amplitude of $\ket{0^n}$ is
$2^{-n}\sum_x(-1)^{f(x)}$: it equals $\pm1$ (probability $1$) when $f$ is constant and
$0$ when $f$ is balanced. The promise ``$f$ is constant or balanced'' is what makes
the one-query, all-zeros test exact.

% ======================================================================
\problem{6 (Grover search, $N=16$)}

\sol
\textbf{(a)} $\sin\tfrac\theta2=\tfrac1{\sqrt{16}}=\tfrac14$, so
$\theta=2\arcsin\tfrac14\approx0.5054$ ($\approx28.96^\circ$). $R_{x_0}=\id-2\proj{x_0}$
reflects about the axis $\ket{\psi_0^\perp}$ and $R_{\psi_0}=2\proj{\psi_0}-\id$
reflects about $\ket{\psi_0}$; a product of two reflections is a rotation by twice the
angle between their mirror axes. Since $\ket{\psi_0}$ sits at angle $\theta/2$ from
$\ket{\psi_0^\perp}$, $G=R_{\psi_0}R_{x_0}$ rotates by $\theta$ toward $\ket{x_0}$.

\textbf{(b)} $p_k=\sin^2\!\bigl((2k+1)\theta/2\bigr)$:
\[
 p_0=0.0625,\ p_1=0.4727,\ p_2=0.9084,\ p_3=0.9613,\ p_4=0.5817,\ p_5=0.1255 .
\]
The optimum is $k^\ast=3$ ($p_3\approx0.961$). At $k=4,5$ the angle
$(2k+1)\theta/2$ passes $\tfrac\pi2$, so the state rotates \emph{past} $\ket{x_0}$ and
the success probability falls back (over-rotation).

\textbf{(c)} For $N=2^n$ with one marked item, $\theta/2=\arcsin(1/\sqrt N)\approx
1/\sqrt N$; maximizing wants $(2k+1)\theta/2\approx\tfrac\pi2$, giving
$k^\ast\approx\tfrac{\pi}{4}\sqrt N$ and query complexity $\bigo(\sqrt N)$.

% ======================================================================
\problem{7 (Hadamard test)}

\sol
\textbf{(a)} After $H$ on the ancilla and controlled-$U$ the state is
$\tfrac1{\sqrt2}(\ket0\ket\psi+\ket1 U\ket\psi)$; a second $H$ gives
$\tfrac12\bigl[\ket0(\ket\psi+U\ket\psi)+\ket1(\ket\psi-U\ket\psi)\bigr]$, so
$p(0)=\tfrac14\norm{\ket\psi+U\ket\psi}^2=\tfrac14(2+2\Re\braket{\psi|U|\psi})
=\tfrac12(1+\Re\braket{\psi|U|\psi})$. Inserting $S\dg=\diag(1,-\ii)$ on the ancilla
between the two Hadamards replaces $\Re$ by $\Im$: $p(0)=\tfrac12(1+\Im\braket{\psi|U|\psi})$.

\textbf{(b)} $U=S=\diag(1,\ii)$, $\ket\psi=\ket+$:
$\braket{+|S|+}=\tfrac12(\braket0+\bra1)(\ket0+\ii\ket1)=\tfrac12(1+\ii)$, so
$\Re=\Im=\tfrac12$. The real-part test predicts $p(0)=\tfrac12(1+\tfrac12)=\tfrac34$ and
the imaginary-part variant also $\tfrac34$, consistent with $\braket{+|S|+}=(1+\ii)/2$.

% ======================================================================
\problem{8 (QFT and phase estimation)}

\sol
\textbf{(a)} For $n=3$,
$\QFT\ket{j}=\tfrac1{\sqrt8}\bigl(\ket0+\ee^{2\pi\ii\,0.j_3}\ket1\bigr)
\bigl(\ket0+\ee^{2\pi\ii\,0.j_2j_3}\ket1\bigr)
\bigl(\ket0+\ee^{2\pi\ii\,0.j_1j_2j_3}\ket1\bigr)$. The circuit uses $3$ Hadamards and
$\binom32=3$ controlled-phase gates ($R_2,R_3$), plus the final bit-reversal swap.

\textbf{(b)} $\varphi=\tfrac38$, $t=3$. After the controlled powers the ancilla is
$\tfrac1{\sqrt8}\sum_{j=0}^{7}\ee^{2\pi\ii(3/8)j}\ket j=\QFT\ket3$, because
$\varphi=k/8$ with $k=3$. Applying $\QFT^{-1}$ returns $\ket3=\ket{011}$, measured with
probability $1$. The readout is \emph{exact} because $\varphi=k/2^t$ is a dyadic
rational with a $t$-bit expansion.

\textbf{(c)} Using $t=n+\lceil\log_2(2+\tfrac1{2\delta})\rceil$ for success
$\ge1-\delta$: with $n=6$, $\delta=0.1$ we get
$\lceil\log_2(2+5)\rceil=\lceil\log_2 7\rceil=3$, hence $t=9$ ancilla qubits.

% ======================================================================
\problem{9 (Ground-state energy estimation)}

\sol
\textbf{(a)} $U\ket{E_j}=\ee^{-\ii H t}\ket{E_j}=\ee^{-\ii E_j t}\ket{E_j}
=\ee^{2\pi\ii\varphi_j}\ket{E_j}$ with $\varphi_j=-\tfrac{E_j t}{2\pi}\bmod1$. QPE
estimates $\varphi_j$; inverting, $E_j=-\tfrac{2\pi}{t}\varphi_j$ (taken in the window
chosen so $E_j t/2\pi\in[0,1)$).

\textbf{(b)} By linearity QPE on $\ket\phi=\sum_j c_j\ket{E_j}$ collapses onto
$\ket{E_j}$ (returning $\varphi_j$) with probability $\abs{c_j}^2$; the ground phase
$\varphi_0$ thus appears with probability $p_0=\abs{c_0}^2=0.4$ per run. For $R$
independent runs, $\Prob(\text{seen at least once})=1-(1-p_0)^R\ge0.99$ requires
$(0.6)^R\le0.01$, i.e.\ $R\ge\ln0.01/\ln0.6\approx9.02$, so $R=10$ runs suffice.

\textbf{(c)} Energy estimation only reads \emph{some} eigenvalue each run and keeps
the smallest observed; it needs the ground state to have nonnegligible overlap, but
\emph{not} a spectral gap, since one is not trying to resolve neighboring levels.
Preparing the ground state (filtering out the excited components) does need a gap to
separate $E_0$ from the rest.

% ======================================================================
\problem{10 (Amplitude estimation)}

\sol
\textbf{(a)} Write $A\ket{0^n}=\sin\tfrac\theta2\ket{\phi_{\mathrm{good}}}
+\cos\tfrac\theta2\ket{\phi_{\mathrm{bad}}}$, so $P=\sin^2\tfrac\theta2$. On the plane
$\operatorname{span}\{\ket{\phi_{\mathrm{good}}},\ket{\phi_{\mathrm{bad}}}\}$ the
iterate $G=R_{\phi_0}R_{\mathrm{good}}$ is a rotation by $\theta$ with eigenvalues
$\ee^{\pm\ii\theta}$; the prepared state is an equal superposition of the two
eigenvectors, so QPE on $G$ returns $\pm\theta$ and hence $P=\sin^2(\theta/2)$.

\textbf{(b)} For $\epsilon=10^{-3}$: amplitude estimation uses
$\bigo(1/\epsilon)\approx10^{3}$ applications of $A$ (and of $G$), whereas classical
Monte~Carlo needs $\bigo(1/\epsilon^2)\approx10^{6}$ samples---a factor $\approx10^3$
fewer queries.

\textbf{(c)} The Hadamard test encodes $\Re\braket{\psi|U|\psi}$ in the bias
$p(0)=\tfrac12(1+\Re\braket{\psi|U|\psi})$. Estimating $p(0)$ by repetition costs
$\bigo(1/\epsilon^2)$ runs (shot noise); running amplitude estimation on the very same
circuit estimates $p(0)$---hence $\Re\braket{\psi|U|\psi}$---to precision $\epsilon$
with $\bigo(1/\epsilon)$ coherent applications. This is the quadratic
(Heisenberg-limited) improvement over plain repetition.

\end{document}
